\section{Произношение согласных звуков}

\begin{longtblr}{
    colspec = { X[1,l,t] X[1,l,t] X[1,l,t] X[3,l,t] },
    rowhead=1,
    row{1} = {font=\bfseries},
}
\toprule
Румынский звук & Русский аналог & Написание & Комментарий \\
\midrule
\textipa{[b p f v m n r z]}
& б, п, ф, в, м, н, р, з
& b, p, f, v, m, n, r, z
& Румынские согласные звуки близки соответствующим русским звукам.
  Отличие в их большей напряженности и отсутствии смягчения перед гласными i, e.
  Смягчения не происходит и перед дифтонгом \textipa{[\t*{ea}]}.
  Звонкие согласные на конце слова не оглушаются.
\\
\textipa{[b' p' f' v' m' n' r' z']}
& пь, бь, фь, вь, мь, нь, рь, зь
& -pi, -bi, -fi, -vi, -mi, -ni, -ri, -zi
& Мягкие звуки близки русским согласным, за которыми на письме следует мягкий знак.
  Румынские звуки произносятся значительно четче и напряженнее, чем соответствующие им русские,
  и в конце слова не оглушаются.
  Мягкие согласные чаще всего встречаются на конце слов.
  Буква i в таких случаях является мягким знаком: она не произносится и ударения не несет.
\\
\textipa{[k]}
& к
& c, ch, k
& Звук \textipa{[k]} передается на письме двояко: ch (перед i, e) либо c (в остальных случаях).
  Буква h в сочетании ch не произносится.
  Буква k используется только в словах иностранного происхождения.
\\
\textipa{[g]}
& г
& g, gh
& Звук \textipa{[g]} передается на письме двояко: gh (перед i, e) либо g (в остальных случаях).
  Буква h в сочетании gh не произносится.
\\
\textipa{[k' g']}
& кь, гь
& chi, ghi
& Звуки \textipa{[k'], [g']} на письме передаются сочетанием букв chi, ghi как в конце слова,
  так и в других позициях (перед произносимым гласным).
\\
\textipa{[kI gI]}
& {ки, ги \\ (не смягч.)}
& {chi, ghi \\ (перед согл.)}
& Сочетания букв chi, ghi перед согласным произносятся как \textipa{[kI], [gI]},
  при этом \textipa{[k], [g]} не смягчаются.
\\
\textipa{[Nk Ng]}
& --
& nc, nch, ng, ngh
& Звук \textipa{[N]} встречается только перед \textipa{[k], [g]}.
  Переход от \textipa{[N]} к следующему согласному \textipa{[k], [g]} происходит плавно.
  Звук \textipa{[N]} появляется и на стыке слов при отсутствии между ними паузы,
  если одно из них оканчивается на -n, а следующее начинается на c, ch, g, gh.
\\
\textipa{[Nk' Ng']}
& --
& -nchi, -nghi
& Сочетания букв nchi, nghi на конце слова звучат как \textipa{[Nk'], [Ng']}.
\\
\textipa{[s t d]}
& с, т, д
& s, t, d
& Произношение звуков \textipa{[s t d]} сходно с соответствующими русскими звуками.
  Не смягчаются перед гласными i, e и перед дифтонгом \textipa{[\t*{ea}]}.
  Звуки \textipa{[s], [d]} могут быть только твердыми.
  Смягчение \textipa{[t]} происходит только в сочетании -ști \textipa{[St']} на конце слова.
\\
\textipa{[h h']}
& х, хь
& h, -hi
& Придыхание в звуке \textipa{[h]} менее заметно, чем в русском х.
  Звук \textipa{[h']} произносится более напряженно, чем русский хь.
  Звуку \textipa{[h']} соответствует написание hi на конце слова при отсутствии ударения.
\\
\textipa{[l l']}
& л, ль
& l, -li
& При произнесении звука \textipa{[l]} следует прижать кончик языка к верхним альвеолам, опустив спинку языка.
  Звуку \textipa{[l']} соответствует сочетание букв li (на конце слова при отсутствии ударения).
\\
\textipa{[S S']}
& ш, шь
& ș, -și
& Звук \textipa{[S]} близок русскому звуку ш.
  Звуку \textipa{[S']} соответствует сочетание букв și (на конце слова при отсутствии ударения).
  В сочетании și (с произносимым i) звук \textipa{[I]} не переходит в ы.
  В сочетании ști, если оно находится в начале слова или на конце слова под ударением, гласный произносится.
  На конце слова в безударном положении сочетание ști произносится как \textipa{[St']}.
\\
\textipa{[Z Z']}
& ж, жь
& j, -ji
& В положении перед согласным или на конце слова звук \textipa{[Z]} не оглушается.
  Звуку \textipa{[Z']} соответствует сочетание букв ji (на конце слова при отсутствии ударения).
\\
\textipa{[ts ts']}
& ц, ць
& ț, -ți
& Звук \textipa{[ts]} в любом положении произносится мягче русского ц.
  Гласный \textipa{[I]} после ț сохраняет свое обычное произношение и не переходит в ы.
  Звуку \textipa{[ts']} соответствует сочетание букв ți (на конце слова при отсутствии ударения).
\\
\textipa{[tS]}
& ч
& ce, ci
& Произношение звука \textipa{[tS]} близко к русскому звуку ч,
  но отличается большей напряженностью и произносится несколько тверже.
  При произнесении \textipa{[tS]} смычка образуется путем прижимания кончика языка к верхним альвеолам.
  Сочетание ci произносится либо как \textipa{[tSI]} --- в середине слова или в конце слова под ударением,
  либо как \textipa{[tS]} --- в конце слова, если i ударения не несет.
  Сочетание cea произносится как \textipa{[tS\t*{ea}]}.
  В сочетаниях cio, ciu буква i не произносится, сочетания звучат как \textipa{[tSo]}, \textipa{[tSu]}.
\\
\textipa{[dZ]}
& --
& ge, gi
& Звук \textipa{[dZ]} можно получить, произнеся максимально слитно и напряженно сочетание букв дж.
  Смычка образуется путем прижимания кончика языка к верхним альвеолам.
  Сочетание gi произносится либо как \textipa{[dZI]} --- в середине слова или в конце слова под ударением,
  либо как \textipa{[dZ]} --- в конце слова, если i ударения не несет.
  Сочетание gea произносится как \textipa{[dZ\t*{ea}]}.
  В сочетаниях gia, gio, giu буква i не произносится,
  сочетания звучат как \textipa{[dZa]}, \textipa{[dZo]}, \textipa{[dZu]}.
  Слова с суффиксом -giu составляют исключение: в них i произносится и несет на себе ударение.
\\
\bottomrule
\end{longtblr}

Примеры:
\begin{description}
    \item[\textipa{[b p f v m n r z]}:] babă, popă, fenomen, vin, iarmaroc, grupă, roză;
        primi \textipa{[prI"mI]}, veac \textipa{[v\t*{ea}k]}, rob \textipa{[rob]}, gaz \textipa{[gaz]}.
    \item[\textipa{[b' p' f' v' m' n' r' z']}:] robi \textipa{[rob']},
        crabi \textipa{[crab']}, \ldots, ursuzi \textipa{[ur"suz']}.
    \item[\textipa{[k g]}:] cheia, chip, cheamă, ghete, ghimp, gheată,
        copac, cracă, goană, greu, târg, kilometru.
    \item[\textipa{[k' g']}:] ochi \textipa{[ok']}, unghi \textipa{[uNg']},
        chiar \textipa{[k'ar]}, ochiuri \textipa{[ok'ur']}, ghiont \textipa{[g'ont]}.
    \item[\textipa{[kI gI]}:] chip \textipa{[kIp]}, ghid \textipa{[gId]}.
    \item[\textipa{[Nk Ng Nk' Ng']}:] încă \textipa{["1Nk@]}, anchetă \textipa{[aN"ket@]}, ungur \textipa{["uNgur]};
        un chinez \textipa{[uNkinez]}, un covor \textipa{[uNkovor]}, un grup \textipa{[uNgrup]};
        unchi \textipa{[uNk']}, unghi \textipa{[uNg']}.
    \item[\textipa{[s]}:] sanie, siberian, se, seamă, seară, sac, casă.
    \item[\textipa{[t d]}:] Dunărea, tătar, tir \textipa{[tIr]}, teafăr \textipa{[t\t*{ea}f@r]},
        dădea \textipa{[d@d\t*{ea}]}; ești \textipa{[\u*{i}eSt']}, vorbești \textipa{[vorbeSt']}.
    \item[\textipa{[h h']}:] hrean, monarhi \textipa{[monarh']}, arhitect \textipa{[arhItekt]}.
    \item[\textipa{[l l']}:] album, calm, mult, canal, copilul, deal, el;
        banal --- banali, egal --- egali, vesel --- veseli.
    \item[\textipa{[S S']}:] șanse, școală, duș, șină, șef, tușea, ști \textipa{[StI]},
        știre \textipa{[StIre]}, crești \textipa{[kreSt']}; cocoș --- cocoși, moș --- moși, frumoși, urși.
    \item[\textipa{[Z Z']}:] joc, jiletcă \textipa{[ZIletk@]}, jar, paj;
        coajă --- coji, paj --- paji, obraji.
    \item[\textipa{[ts ts']}:] lecție, țară, țigan \textipa{[tsI"gan]};
        hoț --- hoți, creț --- creți, stați, vreți.
    \item[\textipa{[tS]}:] face, cine; munci \textipa{[mun"tSI]}, munci \textipa{["muntS]};
        ceas \textipa{[tS\t*{ea}s]}, ciorbă \textipa{[tSorb@]}, ciumă \textipa{[tSum@]}.
    \item[\textipa{[dZ]}:] gem, merge, gigantic, magic; îndrăgi\textipa{[1ndr@dZI]}, dragi\textipa{[dradZ]};
        geam, Caragiale, gioarse, giugiuli;
        giuvaiergiu \textipa{[dZuva\u*{i}er"dZI\u*{u}]}, sacagiu \textipa{[saka"dZI\u*{u}]}.
\end{description}
